%% ma: L11-B16-0060
%% lop: 11
%% chuong: 5
%% bai: 16
%% dang: 11.16.7
%% loai: DS
%% muc: [NB, TH, TH, TH]
%% dap_an: "ĐSSĐ"
%% nguon: "tai-lieu-hoc-tap-mon-toan-lop-11-theo-chuong-trinh-giao-duc-pho-thong-2018.pdf (Lớp Toán Thầy Duy, Phạm Lê Duy, THPT Chu Văn An), Chương 3 Bài 2, Phần 2, Câu 3 (tr.377)"
%% kien_thuc: "Lớp 11 Bài 16 (giới hạn của hàm số)"
%% trang_thai: cho-duyet
%% ly_do: "Dạng mới đề xuất (11.16.x), chờ giáo viên duyệt dạng; tài liệu gốc không có lời giải mẫu, đáp án do kho tự giải và kiểm bằng sympy."
%% ghi_chu: "Đề gốc in ý c), d) là x->1^- và x->1^+; đáp án gốc (S, Đ) chỉ khớp đồ thị nếu là x->(-1)^- và x->(-1)^+ (mất dấu trừ), nên kho đã sửa thành (-1). Hình TikZ do kho dựng lại theo đồ thị gốc (dạng gần đúng: tiệm cận đứng x=-1, tiệm cận ngang y=2, đồ thị qua (0;2) và (1;0)); đồ thị gốc không cho dáng điệu chính xác ở vô cực nên ý a), b) dựa vào tiệm cận ngang y=2 vẽ trên hình."
\begin{ex}
Cho đồ thị hàm số $y=f(x)$ như hình vẽ (đồ thị có tiệm cận đứng $x=-1$ và tiệm cận ngang $y=2$).
\begin{center}
\begin{tikzpicture}[x=0.7cm,y=0.5cm]
\draw[truc] (-5.4,0)--(4.4,0) node[below]{$x$};
\draw[truc] (0,-1)--(0,7.3) node[left]{$y$};
\draw[blue!60!black] (-1,-1)--(-1,7.3);
\draw[blue!60!black] (-5.4,2)--(4.4,2);
\draw[cyan!70!blue,thick,domain=-5:-1.8,samples=60] plot (\x,{2+4/(-\x-1)});
\draw[cyan!70!blue,thick,domain=-0.55:1,samples=60] plot (\x,{2*(1-\x)/(1+\x)});
\draw[cyan!70!blue,thick,domain=1:4.2,samples=40] plot (\x,{2-2/\x});
\path (0,2) node[above right]{\footnotesize$2$} (-1,0) node[below left]{\footnotesize$-1$} (0,0) node[below left]{\footnotesize$O$} (1,0) node[below]{\footnotesize$1$};
\end{tikzpicture}
\end{center}
\choiceTF
{\True $\lim\limits_{x\to+\infty}f(x)=2$.}
{$\lim\limits_{x\to-\infty}f(x)=-\infty$.}
{$\lim\limits_{x\to(-1)^-}f(x)=2$.}
{\True $\lim\limits_{x\to(-1)^+}f(x)=+\infty$.}
\loigiai{\sloppy Từ đồ thị: đường $y=2$ là tiệm cận ngang nên $f(x)\to2$ khi $x\to\pm\infty$; đường $x=-1$ là tiệm cận đứng, đồ thị đi lên $+\infty$ ở cả hai phía của $x=-1$.
a) $\lim\limits_{x\to+\infty}f(x)=2$. Đúng.
b) $\lim\limits_{x\to-\infty}f(x)=2\ne-\infty$. Sai.
c) $\lim\limits_{x\to(-1)^-}f(x)=+\infty\ne2$. Sai.
d) $\lim\limits_{x\to(-1)^+}f(x)=+\infty$. Đúng.}
\end{ex}
